% ============================================================
% RESULTADOS
% ============================================================

% -------------------------------------------------------
\section{Resultados}
% -------------------------------------------------------

\begin{frame}[plain]
    \vfill
    \centering
    \begin{beamercolorbox}[sep=12pt,center,shadow=false,rounded=true]{title}
        {\Large\bfseries Resultados}\\[6pt]
        {\normalsize Evidência empírica e decomposição do bem-estar}
    \end{beamercolorbox}
    \vfill
\end{frame}

\begin{frame}{Decomposição do bem-estar}
    \textbf{Contrafactual principal:} rede rodoviária de Brasília \textit{vs.}\ ausência de estradas.

    \vspace{0.25cm}

    \begin{columns}[T]
        \begin{column}{0.48\textwidth}
            \textbf{Ganho total}
            \begin{itemize}
                \item Bem-estar médio $+2{,}8\%$.
                \item Custos de comércio $\downarrow 9\%$; custos de migração $\downarrow 8\%$.
            \end{itemize}

            \vspace{0.2cm}

            \textbf{Decomposição dos canais}
            \begin{itemize}
                \item \textbf{Comércio:} $76\%$ do ganho ($+2{,}1$ p.p.).
                \item \textbf{Migração:} $24\%$ do ganho ($+0{,}7$ p.p.).
            \end{itemize}
        \end{column}
        \begin{column}{0.48\textwidth}
            \textbf{Intuição}
            \begin{itemize}
                \item Estradas $\downarrow$ $\tau_{odt}$ $\Rightarrow$ $P_{dt}$ $\downarrow$ $\Rightarrow$ $V_{dt}$ $\uparrow$: ganho via preços.
                \item Estradas $\downarrow$ $\kappa_{odt}$ $\Rightarrow$ realocação de trabalho: ganho via migração.
                \item Canal comercial domina porque o $\theta = 4$ amplifica reduções de custo iceberg.
                \item Comércio e migração aumentaram $\approx 20\%$ cada.
            \end{itemize}
        \end{column}
    \end{columns}

    \vspace{0.2cm}
\end{frame}

\begin{frame}{Heterogeneidade espacial}
    \begin{columns}[T]
        \begin{column}{0.52\textwidth}
            \textbf{Com migração custosa (cenário base)}
            \begin{itemize}
                \item Ganhos variam de \textbf{1\% a 15\%} entre estados.
                \item Estados mais conectados à rede MST ganham desproporcionalmente.
                \item Fricções de migração impedem equalização rápida de utilidades.
            \end{itemize}

            \vspace{0.3cm}

            \textbf{Com mobilidade perfeita}
            \begin{itemize}
                \item Ganhos se equalizam entre todos os estados.
                \item Mostra que custos de migração são centrais para gerar heterogeneidade espacial.
            \end{itemize}
        \end{column}
        \begin{column}{0.44\textwidth}
            \textbf{Contrafactual: Rio de Janeiro}
            \begin{itemize}
                \item Rede alternativa centrada no Rio substituiria a rede de Brasília.
                \item Bem-estar $\approx 0{,}2\%$ \textbf{menor} com a rede do Rio.
                \item Resultado: a localização geográfica de Brasília foi, em si, vantajosa para a integração nacional.
            \end{itemize}
        \end{column}
    \end{columns}

    \vspace{0.15cm}
\end{frame}
